\documentclass[10pt,a4paper]{amsart}
\usepackage[margin=19mm]{geometry}
\usepackage{amsmath,amssymb,mathtools}
\usepackage[scr=rsfs,cal=euler]{mathalfa}
\usepackage{xfrac}
\usepackage[unicode,hidelinks,bookmarks=false]{hyperref}
\newcommand{\ie}{i.e.\ }
\newcommand{\cf}{cf.\ }
\newcommand{\resp}{respectively\ }
\newcommand{\Subsection}{\subsection}
\providecommand{\nobreakdash}{\nobreak\hbox{-}}
\setlength{\emergencystretch}{2em}
\multlinegap0pt
\usepackage{amssymb}
\usepackage{tcolorbox}
\usepackage{tikz}
\usepackage{xcolor}
\usepackage{mathtools}
\newtheorem{proposition}{Proposition}
\newcommand{\R}{\mathbb R}
\renewcommand{\ge}{\geqslant}
\renewcommand{\le}{\leqslant}
\title{A High-Frequency Uncertainty Principle for\\ the Fourier-Bessel Transform}
\author{Benjamin Jaye, Rahul Sethi}
\date{Source: arXiv:2509.25500v2 (CC BY 4.0)}
\begin{document}
\maketitle

\noindent\emph{Source and licence.} A High-Frequency Uncertainty Principle for\\ the Fourier-Bessel Transform. Version arXiv:2509.25500v2 (CC BY 4.0), distributed by arXiv under the Creative Commons Attribution 4.0 International licence, http://creativecommons.org/licenses/by/4.0/, as recorded in the arXiv licence metadata for this exact version and shown on the abstract page https://arxiv.org/abs/2509.25500v2. Full licence notice: \texttt{licenses/arxiv-2509.25500-CC-BY-4.0.txt}.\par\smallskip

\noindent\textit{Editorial scope.} Counted unit: the article's proposition relative-density-equivalence (source0.tex lines 258--261). The article's complete original proof of it is delivered with it (source0.tex lines 262--275, one \texttt{proof} environment of 14 lines). No mathematics is written, repaired or re-derived: the statement and the proof are the article's own lines, in the article's own notation, and no retained line is substituted or re-flowed.  No further source-local context is needed: every cross-reference of every retained line resolves inside the two delivered spans.  Nothing was omitted from any retained span, so the package carries no omission record.  No retained line contains a citation, so the wrapper carries no bibliography and no bibliography entry.  No retained line contains a figure environment, an image-inclusion command or drawing code, so no image is delivered and the package declares no asset dependency.  Editorial formatting only, in addition: an AMS \texttt{amsart} preamble that restates the article's own declarations and macros used by the retained lines, namely the environments \texttt{proposition} on separate counters, together with the standard packages those lines need.  The retained environments print in this document as Proposition 1 in order of appearance, whereas the article prints the same items with the numbering of its own full document.  The complete native source of the article as distributed in the arXiv e-print for this exact version is bundled unchanged as \texttt{sources/186018/source0.tex}.
\begin{proposition}
\label{relative-density-equivalence}
    Let $\alpha > -1/2$. A set $E \subset \R^+$ is $\mu_\alpha$-relatively dense if and only if it is Lebesgue relatively dense. 
\end{proposition}

\begin{proof}
    Suppose $\gamma > 0$ is such that $|E \cap [r, r+1]| \ge \gamma$ for all $r \ge 0$. We will find $\widetilde{\gamma} > 0$ such that $\mu_\alpha(E \cap [r,r+1]) \ge \widetilde{\gamma} \mu_\alpha([r,r+1])$. Using $r^{2\alpha + 1} \le x^{2\alpha + 1} \le (r+1)^{2\alpha + 1},$ we have the inequalities
    $$c_1 r^{2\alpha + 1} |E \cap [r,r+1]| \le \mu_\alpha(E \cap [r,r+1]) \le c_1 (r+1)^{2\alpha + 1}\, |E \cap [r,r+1]|,$$
    and $$c_1 r^{2\alpha + 1} \le \mu_\alpha([r,r+1]) \le c_1 (r+1)^{2\alpha + 1},$$
    where $c_1 := \frac{2\pi^{\alpha + 1}}{\Gamma(\alpha + 1)}$. Combining the two, we get
    $$\frac{\mu_\alpha(E \cap [r,r+1])}{\mu_\alpha([r,r+1])} \ge \gamma \left(\frac{r}{r+1}\right)^{2\alpha + 1}  \ge \gamma \left(\frac{\frac{\gamma}{2}}{\frac{\gamma}{2}+1}\right)^{2\alpha + 1} > 0,$$
    when $r \ge \frac{\gamma}{2}$. If $r \le \frac{\gamma}{2}$, we split $I := [r,r+1]$ as $I = J \cup J'$, where $J = [r + \frac{\gamma}{2}, r + 1]$ and $J' = [r, r+ \frac{\gamma}{2}]$. We have $$\mu_\alpha(E \cap I) \ge \mu_\alpha(E \cap J) = c_1\int_{E \cap J} x^{2\alpha + 1}\, dx \ge c_1 \left(\frac{\gamma}{2} \right)^{2\alpha + 1} |E \cap J| \ge c_1 \left(\frac{\gamma}{2} \right)^{2\alpha + 2},$$
    as
    $$|E \cap J| = |E \cap I| - |E \cap J'| \ge \gamma - \frac{\gamma}{2} = \frac{\gamma}{2}.$$
    On the other hand, $$\mu_\alpha(I) = c_1\int_I x^{2\alpha + 1}\, dx \le c_1 (r+1)^{2\alpha + 1} \le c_1 \left(\frac{\gamma}{2} + 1 \right)^{2\alpha + 1},$$
    giving 
    $$\frac{\mu_\alpha(E \cap [r,r+1])}{\mu_\alpha([r,r+1])} \ge \frac{\gamma}{2} \left(\frac{\frac{\gamma}{2}}{\frac{\gamma}{2}+1}\right)^{2\alpha + 1},$$
    for all $r \le \frac{\gamma}{2}$. We obtain $\mu_\alpha(E \cap [r,r+1]) \ge \widetilde{\gamma} \mu_\alpha([r,r+1])$ with $\widetilde{\gamma} := \frac{\gamma}{2} \left(\frac{\frac{\gamma}{2}}{\frac{\gamma}{2}+1}\right)^{2\alpha + 1}$. The reverse implication is entirely analogous.
\end{proof}

\end{document}
